%9. Conclusion:

%===============================================================================================
% - real-world validation still needed to disparage any remaining doubt about acceptance stated in surveys: Swiss referenda on increasing foreign aid confirm survey results
%===============================================================================================

Despite the extensive survey evidence reviewed throughout this chapter, real-world validation remains limited. Assessing the external validity of these findings requires examining whether stated support persists when citizens make binding political choices with actual budgetary consequences. The 2006 Swiss referendum on the Federal Act on Cooperation with the States of Eastern Europe provides relevant evidence. With a turnout of 44.98\%, 53.4\% of voters approved the Act, which aimed, in part, to provide financial assistance to Eastern European countries and established the legal basis for Switzerland's contribution to reducing economic and social disparities within the enlarged European Union.\footnote{The official results of the referendum are available at \url{https://www.bk.admin.ch/ch/f/pore/va/20061126/det526.html}. For a presentation of the Act and its role as the legal basis for Switzerland's enlargement contribution, see \url{https://www.bk.admin.ch/bk/fr/home/documentation/votations/votation-populaire-du-26-novembre-2006.html}.} This result shows that a majority of voters accepted a binding policy involving cross-border financial assistance. One referendum, however, cannot establish that the support observed in surveys generalizes across policies, countries, and electoral contexts.

%===============================================================================================
% - given these results, even though public demand may not be organized, it's very unlikely that we'd see protests or popular rejection of int'l redistr policies
%===============================================================================================

As discussed above, respondents generally support international redistribution policies when surveys present them with concrete proposals, yet these policies rarely emerge spontaneously when respondents identify their current concerns \citep{fabre_majority_2025,fabre_public_2026}. Pluralistic ignorance, underestimation of public support, and low political salience may account for this discrepancy \citep{andre_globally_2024,fabre_public_2026,mildenberger_beliefs_2017}. The evidence therefore points to latent acceptance rather than an organized political demand. It also provides little basis for expecting broad popular opposition to international redistribution. The central political challenge lies in converting diffuse and low-salience support into sustained collective mobilization.

%===============================================================================================
% - more research needed to understand elite attitudes on these issues
%===============================================================================================

Public opposition alone is therefore unlikely to fully explain the limited policy action in this area. Attention should turn to the actors who select the issues and policy proposals that enter public debate, including political elites, elected representatives, political parties, interest groups, and the media. Future research should examine their attitudes toward international redistribution, their perceptions of public opinion and electoral risks, their assessments of administrative feasibility, and their exposure to organized interests. Existing research shows that political representatives and other elites can systematically misperceive public preferences, often underestimating support for progressive or climate-related policies \citep{mildenberger_beliefs_2017,broockman_bias_2018,pilet_politicians_2023}. Such evidence would help distinguish among genuine elite opposition, misperceptions of public constraints, concerns about implementation, and the institutional weakness of groups advocating international redistribution.

Overall, the evidence reviewed in this chapter shifts the burden of proof. Explanations for the political absence of international redistribution must now account for its low salience, the weakness of organized mobilization, and the choices and beliefs of political elites, rather than assume widespread public opposition.

